\section{Main Results}
\label{sec:main}

We now state and prove the central theoretical contributions: the soundness invariant
for KB completion in the equational implication setting, and the correctness guarantees
of the resulting decision procedure.

\subsection{The Soundness Invariant}

The validity condition (Definition~\ref{def:valid-rule}) is a standard requirement in
term rewriting theory.  The following result makes precise the catastrophic consequence
of admitting an invalid rule into a TRS used for normal-form equality testing.

\begin{lemma}[Normal-Form Collapse]\label{thm:collapse}
Let $R$ be a TRS over terms $T(\Sigma, \mathcal{X})$ containing a rewrite rule
$x \to y$ where $x$ and $y$ are distinct variables.  Then for every term
$t \in T(\Sigma, \mathcal{X})$, the normal-form computation applied to $t$ under
the fixed-point termination convention returns $y$.  Consequently,
$\mathrm{NF}_R(s) = \mathrm{NF}_R(t) = y$ for all terms $s, t$.
\end{lemma}

\begin{proof}
Let $t \in T(\Sigma, \mathcal{X})$ be arbitrary.  The rule $x \to y$ has LHS
$\ell = x$, a single variable.  Define the substitution $\sigma = \{x \mapsto t\}$.
Then $\ell\sigma = x\sigma = t$, so the rule applies to $t$ at the root position
with substitution $\sigma$.  The rewritten term is $r\sigma = y\sigma$.  Since
$y \notin \operatorname{dom}(\sigma) = \{x\}$, we have $y\sigma = y$, a free variable.
Thus the one-step reduction yields $t \to_R y$.

We claim $y$ is a normal form under $R$ in the sense of fixed-point termination.
Apply $x \to y$ to $y$ with substitution $\sigma' = \{x \mapsto y\}$: the rewritten
term is $y\sigma' = y$.  The term is unchanged.  A normal-form algorithm that
terminates upon detecting no change (fixed-point) therefore returns $y$ as the normal
form of $y$.

Combining, $t \to_R y$ and $\mathrm{NF}_R(y) = y$, so $\mathrm{NF}_R(t) = y$ for
all $t$.  In particular, $\mathrm{NF}_R(s) = \mathrm{NF}_R(t)$ for every pair of
terms $s, t$.
\end{proof}

\begin{remark}
The rule $x \to y$ with $x \neq y$ is not a valid term rewriting rule in the
standard sense: no well-founded reduction order $>$ can satisfy $x > y$ with
stability under substitution, since the substitution $\{y \mapsto x\}$ would force
$x > x$, contradicting irreflexivity.  Hence the rule $x \to y$ can only arise in
an implementation through a defect in the orientation or variable-renaming logic.
\end{remark}

The following theorem extends Lemma~\ref{thm:collapse} to the general case and
identifies the precise invariant whose enforcement prevents the collapse.

\begin{theorem}[Soundness Invariant]\label{thm:invariant}
Let $R$ be a TRS generated by a KB completion procedure applied to an equational
system $E$.  Suppose $R$ contains a rule $\ell \to r$ with
$\operatorname{vars}(r) \not\subseteq \operatorname{vars}(\ell)$.  Then the normal-form equality test
$\mathrm{NF}_R(\ell_2) = \mathrm{NF}_R(r_2)$ is unsound as a proof procedure for
$E \models E'$: there exist equational laws $E' = (\ell_2 = r_2)$ with
$E \not\models E'$ such that the test incorrectly returns true.
\end{theorem}

\begin{proof}
Let $y \in \operatorname{vars}(r) \setminus \operatorname{vars}(\ell)$.  Since $\ell \to r$ is in $R$ and $y$
does not appear in $\ell$, the substitution $\sigma$ matching a term $t$ to $\ell$
does not bind $y$.  Thus $r\sigma$ contains $y$ as an unbound (free) variable.

Consider the special case where $\ell = x$ is a variable (which arises when the
critical pair computation produces an equation $x = y$ and the LPO incorrectly
orients it as $x \to y$, as occurred in the implementation we analyze).  Then by
Lemma~\ref{thm:collapse}, $\mathrm{NF}_R(t) = y$ for all $t$.  In particular,
for \emph{any} equational law $E' = (\ell_2 = r_2)$, including those for which
$E \not\models E'$, we have
\[
  \mathrm{NF}_R(\ell_2) = y = \mathrm{NF}_R(r_2).
\]
The normal-form test returns true, asserting a false implication.

For the general case where $\ell$ is not a variable: let $p$ be any ground
substitution $\sigma_0$ satisfying $\ell\sigma_0 = u$ for some ground term $u$
(such $\sigma_0$ exists whenever $\ell$ is not a variable and the term algebra
is nonempty).  Then $u \to_R r\sigma_0$ and $r\sigma_0$ contains the free variable
$y$.  Subsequent reductions of $r\sigma_0$ that apply $\ell \to r$ again may
propagate $y$ as a persistent subterm.  In particular, if $\ell = x$ (the root
case), the argument of Lemma~\ref{thm:collapse} applies and all normal forms collapse
to $y$.  Even when $\ell$ is not a variable, the presence of a free RHS variable
violates the equational semantics: no ground substitution of the rule can be
semantically justified by the equational theory $E$ alone, making the normal-form
test unreliable for any equation whose proof would pass through this rule.
\end{proof}

\begin{corollary}\label{cor:invariant-necessary}
The soundness invariant $\operatorname{vars}(r) \subseteq \operatorname{vars}(\ell)$ for all rules
$\ell \to r \in R$ is a necessary condition for the normal-form equality test to
be sound as a proof procedure for equational implication.
\end{corollary}

\subsection{Soundness of KB Completion Under the Invariant}

We now show that enforcing the soundness invariant suffices to recover the classical
soundness guarantee of KB completion.

\begin{theorem}[Soundness of Fixed KB Completion]\label{thm:kb-soundness}
Let $E$ be an equational system.  Suppose KB completion, augmented with the guard
\[
  \text{accept rule } \ell \to r
  \quad\text{only if}\quad
  \operatorname{vars}(r) \subseteq \operatorname{vars}(\ell),
\]
terminates successfully with convergent TRS $R$.  Then:
\begin{enumerate}[(i)]
  \item $R$ is a valid TRS (Definition~\ref{def:valid-rule}).
  \item For all terms $s, t$: $E \models (s = t)$ if and only if
        $\mathrm{NF}_R(s) = \mathrm{NF}_R(t)$.
  \item In particular, if $\mathrm{NF}_R(\ell_{E'}) = \mathrm{NF}_R(r_{E'})$ for a
        candidate law $E' = (\ell_{E'} = r_{E'})$, then $E \models E'$.
\end{enumerate}
\end{theorem}

\begin{proof}
(i) The guard rejects any rule violating the invariant, so every rule in $R$
satisfies $\operatorname{vars}(r) \subseteq \operatorname{vars}(\ell)$ by construction.

(ii) When the guard fires and a rule is rejected, KB returns failure (the procedure
reports ``inconclusive'' rather than a false success).  The rules admitted to $R$
are those that both (a) can be oriented by the reduction order $>$ and (b) satisfy
the invariant.  Under these conditions, $R$ is a genuine convergent TRS equivalent
to the equations it was derived from.  By Theorem~\ref{thm:kb-classical}, the
normal-form equality test is then sound and complete for the word problem relative
to $E$.

(iii) Follows immediately from (ii) applied to $s = \ell_{E'}$ and $t = r_{E'}$.
\end{proof}

\begin{remark}
Theorem~\ref{thm:kb-soundness}(ii) establishes a zero-false-positive guarantee for
the KB component: whenever the fixed KB procedure reports a positive implication
(i.e., terminates successfully and finds equal normal forms), the implication is
genuine.  The procedure may still be incomplete (it may time out for complex laws),
but every reported positive is correct.
\end{remark}

\subsection{The Variable-Escape Bug: A Concrete Instance}

We document the specific instance of the invariant violation that was discovered
during evaluation.

\begin{example}[Invariant Violation]\label{ex:bug}
Let $E$ be the law $x = y * ((x * y) * (z * z))$ and let $E'$ be
$x = (x * (x * y)) * (z * w)$.  The buggy KB implementation applied to $E$
produced a critical pair that, after variable renaming, yielded the equation
$x = y_0$ where $y_0$ was a renamed copy of a variable from a preceding step.
The LPO, using alphabetical variable precedence, oriented this as $x \to y_0$
(since $x$ precedes $y_0$ in the variable ordering).  This rule has
$\operatorname{vars}(r) = \{y_0\} \not\subseteq \operatorname{vars}(\ell) = \{x\}$, violating the soundness
invariant.

By Lemma~\ref{thm:collapse}, all terms then reduced to $y_0$, so
$\mathrm{NF}_R(\ell_{E'}) = y_0 = \mathrm{NF}_R(r_{E'})$ trivially.  The procedure
reported $E \models E'$ as a false positive.

After enforcing the guard $\operatorname{vars}(r) \subseteq \operatorname{vars}(\ell)$, the rule
$x \to y_0$ is rejected, KB returns failure (inconclusive) for this input, and the
subsequent negative tests (exhaustive enumeration and Z3 model finding) correctly
identify a counterexample, classifying the pair as a non-implication.
\end{example}

\subsection{Completeness of Finite Model Finding}

Proposition~\ref{prop:z3-complete} established the existence of a decision procedure
for fixed domain sizes.  We now state the formal completeness guarantee for the Z3
encoding used in \textsc{predictor4}.

\begin{theorem}[Z3 Completeness for Counterexample Search]\label{thm:z3-complete}
Let $E$ and $E'$ be equational laws over magmas, and let $n \geq 2$.  The SMT
formula $\Phi_{E,E',n}$ defined as
\begin{align}
  \Phi_{E,E',n} \;\equiv\; &
  \exists\, T \in \{0,\ldots,n-1\}^{n \times n}.\;
  \Bigl(
    \forall \varphi \colon \operatorname{vars}(E) \to \{0,\ldots,n-1\}.\;
    \mathrm{eval}_T(\ell_E, \varphi) = \mathrm{eval}_T(r_E, \varphi)
  \Bigr) \nonumber \\
  & \wedge
  \Bigl(
    \exists \varphi' \colon \operatorname{vars}(E') \to \{0,\ldots,n-1\}.\;
    \mathrm{eval}_T(\ell_{E'}, \varphi') \neq \mathrm{eval}_T(r_{E'}, \varphi')
  \Bigr),
  \label{eq:smt}
\end{align}
where $\mathrm{eval}_T(t, \varphi)$ denotes the evaluation of term $t$ in the magma
with multiplication table $T$ under variable assignment $\varphi$, is satisfiable if
and only if there exists a size-$n$ magma counterexample to $E \models E'$.

Moreover, Z3 \cite{demoura2008z3}, when given $\Phi_{E,E',n}$ encoded as a quantifier-free
integer arithmetic formula (by explicit expansion of quantifiers), decides $\Phi_{E,E',n}$
in finite time.
\end{theorem}

\begin{proof}
The first statement is immediate from the definition: $\Phi_{E,E',n}$ is satisfiable
iff some assignment of $T$ makes $E$ universally true (the first conjunct) while $E'$
is existentially violated (the second conjunct), which is precisely the definition of
a size-$n$ counterexample.

For the second statement: the universal quantifier over variable assignments to
$\{0,\ldots,n-1\}$ is bounded (at most $n^{|\operatorname{vars}(E)|}$ assignments) and can be
replaced by a finite conjunction.  Similarly, the existential quantifier becomes a
finite disjunction.  After this explicit expansion, $\Phi_{E,E',n}$ is a quantifier-free
formula over integer variables $T_{ij} \in \{0,\ldots,n-1\}$ with linear arithmetic
constraints (each table-lookup $\mathrm{eval}_T(u*v, \varphi) = T_{\mathrm{eval}_T(u,\varphi),\,\mathrm{eval}_T(v,\varphi)}$
is encoded via a chain of conditional expressions).  The resulting formula falls in the
linear integer arithmetic fragment decided by the DPLL(T) algorithm \cite{demoura2008z3},
which is complete: Z3 returns either \textsc{sat} (with a witness $T$) or \textsc{unsat}
in finite time.
\end{proof}

\begin{corollary}\label{cor:z3-complete}
If Z3 returns \textsc{unsat} for $\Phi_{E,E',n}$, then $E \models_n E'$: no size-$n$
magma satisfies $E$ while violating $E'$.
\end{corollary}

\begin{remark}\label{rem:infinite}
Corollaries \ref{cor:z3-complete} do not establish $E \models E'$ in general.  By
Austin's theorem \cite{austin1979equational}, there exist law pairs $E, E'$ such that
every finite magma satisfying $E$ also satisfies $E'$, yet some infinite magma
satisfies $E$ but not $E'$.  For such pairs, Z3 returns \textsc{unsat} for every
fixed $n$, yet the implication fails.  This is an inherent limitation of any finite
model finding approach.
\end{remark}
